\begin{abstract}
Moroccan agriculture faces persistent challenges related to crop diseases, water scarcity, climate variability, and limited access to timely agronomic expertise, particularly for small and medium-scale farmers. In this context, digital technologies and artificial intelligence offer promising opportunities to improve the accessibility, responsiveness, and relevance of agricultural advisory services. This project proposes \textbf{AgriConnect}, an intelligent agricultural assistance platform designed to support Moroccan farmers through a mobile-first, multimodal, and multi-agent architecture. The system enables farmers to interact through WhatsApp using text messages and crop images, while a central orchestrator coordinates several specialized agents dedicated to conversational support, image analysis, knowledge retrieval, weather-aware reasoning, contextual personalization, and expert escalation. To improve reliability, the platform combines retrieval-augmented generation with a curated agricultural knowledge base and integrates local weather information to contextualize recommendations. In situations where uncertainty is high or the case is critical, the system supports a human-in-the-loop approach by escalating the interaction to agronomists through a dedicated web interface. The project therefore aims to provide a more accessible, grounded, and context-aware form of agricultural assistance adapted to the Moroccan context. This report presents the design, architecture, implementation, and practical relevance of AgriConnect as an applied intelligent support system for digital agriculture.
\end{abstract}